% Options for packages loaded elsewhere
\PassOptionsToPackage{unicode}{hyperref}
\PassOptionsToPackage{hyphens}{url}
\documentclass[
  ignorenonframetext,
]{beamer}
\newif\ifbibliography
\usepackage{pgfpages}
\setbeamertemplate{caption}[numbered]
\setbeamertemplate{caption label separator}{: }
\setbeamercolor{caption name}{fg=normal text.fg}
\beamertemplatenavigationsymbolsempty
% remove section numbering
\setbeamertemplate{part page}{
  \centering
  \begin{beamercolorbox}[sep=16pt,center]{part title}
    \usebeamerfont{part title}\insertpart\par
  \end{beamercolorbox}
}
\setbeamertemplate{section page}{
  \centering
  \begin{beamercolorbox}[sep=12pt,center]{section title}
    \usebeamerfont{section title}\insertsection\par
  \end{beamercolorbox}
}
\setbeamertemplate{subsection page}{
  \centering
  \begin{beamercolorbox}[sep=8pt,center]{subsection title}
    \usebeamerfont{subsection title}\insertsubsection\par
  \end{beamercolorbox}
}
% Prevent slide breaks in the middle of a paragraph
\widowpenalties 1 10000
\raggedbottom
\AtBeginPart{
  \frame{\partpage}
}
\AtBeginSection{
  \ifbibliography
  \else
    \frame{\sectionpage}
  \fi
}
\AtBeginSubsection{
  \frame{\subsectionpage}
}
\usepackage{iftex}
\ifPDFTeX
  \usepackage[T1]{fontenc}
  \usepackage[utf8]{inputenc}
  \usepackage{textcomp} % provide euro and other symbols
\else % if luatex or xetex
  \usepackage{unicode-math} % this also loads fontspec
  \defaultfontfeatures{Scale=MatchLowercase}
  \defaultfontfeatures[\rmfamily]{Ligatures=TeX,Scale=1}
\fi
\usepackage{lmodern}
\ifPDFTeX\else
  % xetex/luatex font selection
\fi
% Use upquote if available, for straight quotes in verbatim environments
\IfFileExists{upquote.sty}{\usepackage{upquote}}{}
\IfFileExists{microtype.sty}{% use microtype if available
  \usepackage[]{microtype}
  \UseMicrotypeSet[protrusion]{basicmath} % disable protrusion for tt fonts
}{}
\makeatletter
\@ifundefined{KOMAClassName}{% if non-KOMA class
  \IfFileExists{parskip.sty}{%
    \usepackage{parskip}
  }{% else
    \setlength{\parindent}{0pt}
    \setlength{\parskip}{6pt plus 2pt minus 1pt}}
}{% if KOMA class
  \KOMAoptions{parskip=half}}
\makeatother
\setlength{\emergencystretch}{3em} % prevent overfull lines
\providecommand{\tightlist}{%
  \setlength{\itemsep}{0pt}\setlength{\parskip}{0pt}}
\usepackage{bookmark}
\IfFileExists{xurl.sty}{\usepackage{xurl}}{} % add URL line breaks if available
\urlstyle{same}
\hypersetup{
  hidelinks,
  pdfcreator={LaTeX via pandoc}}

\author{\texorpdfstring{}{}}
\date{}

\begin{document}

\section{Scope and related work: landscape and tool
categories}\label{scope-and-related-work-landscape-and-tool-categories}

\begin{frame}{COGANT in the program-analysis landscape}
\protect\phantomsection\label{cogant-in-the-program-analysis-landscape}
COGANT sits at the intersection of three established research areas:
machine learning for source code, graph-based program representations,
and active-inference behavioral modeling. The following subsections
place it against each, emphasizing what COGANT provides as
infrastructure rather than as a novel model or benchmark.

\textbf{Machine learning for big code} surveys cover naturalness,
representation learning, and task taxonomy {[}@allamanis2018survey{]}.
COGANT contributes \textbf{infrastructure}: a documented IR, pipeline,
and export contract rather than a new benchmark or model architecture.

\textbf{Graph neural networks} unify message-passing frameworks for
relational data {[}@wu2020comprehensive; @scarselli2009graph{]}.
Although COGANT's primary output is the Active Inference Institute's
Generalized Notation Notation, its program graph can be materialised as
optional tensor views compatible with these frameworks (PyG, DGL) so
that graph-neural-network model papers can cite a specific tensor
layout.

\textbf{Code property graphs} and related security-oriented graph
constructions show how to merge syntactic and semantic edges for
analysis {[}@yamaguchi2014modeling{]}. COGANT's program graph is cognate
but emphasizes \textbf{ML export}, confidence scoring, and multi-stage
IR refinement.
\end{frame}

\begin{frame}{Related tool categories}
\protect\phantomsection\label{related-tool-categories}
\begin{itemize}
\tightlist
\item
  \textbf{Compiler and LLVM IRs} --- richer semantics, heavier
  toolchain; COGANT favors lightweight extraction for dataset building.
\item
  \textbf{SCA and linters} --- rule enforcement focus; COGANT focuses on
  graph generation for downstream learning.
\item
  \textbf{Neural code models} (code2vec, CodeBERT, etc.) --- often
  consume tokens or AST paths; COGANT supplies \textbf{explicit program
  graphs} that graph-neural-network-centric research can ingest
  directly.
\end{itemize}
\end{frame}

\end{document}
